% Apuntes de teoría en PDF · 4.º ESO · Unidad 3: Polinomios (versión en castellano)
% Mateix estil que els apunts de Batxillerat (franja i caixes de teoria en blau, exemples en roig), però més visuals,
% curts i amb dibuixos, per a l'alumnat de 3r ESO. Es compila amb pdflatex dues vegades (per als enllaços de l'índex),
% amb logo-fentmates.pdf (_eines/logo/pdf/horitzontal-clar.pdf) al costat.
\documentclass[11pt,a4paper]{article}
\usepackage[utf8]{inputenc}
\usepackage[T1]{fontenc}
\usepackage[spanish,es-noshorthands,es-nodecimaldot,es-nolists]{babel}
\usepackage[sfdefault,book,lining]{FiraSans}
\usepackage[italic,noplusnominus,noequal]{mathastext}
\linespread{1.18}
\usepackage[a4paper,top=2.4cm,bottom=2.2cm,left=1.6cm,right=1.6cm,headheight=1.1cm,headsep=0.5cm,footskip=1cm]{geometry}
\usepackage{amsmath,amssymb,cancel}
\usepackage{icomma}
\usepackage{graphicx}
\usepackage[table]{xcolor}
\usepackage{tikz}\usetikzlibrary{calc,positioning,arrows.meta,decorations.pathreplacing,shapes.geometric}
\usepackage{fontawesome5}
\usepackage{ccicons}
\usepackage{fancyhdr}
\usepackage{enumitem}
\usepackage{array}
\usepackage{needspace}
\usepackage[most]{tcolorbox}
\usepackage{titlesec}
\usepackage[hidelinks]{hyperref}
\hypersetup{pdftitle={Apuntes de teoría. Unidad 3: Polinomios. 4.º ESO},pdfsubject={Matemáticas A 4.º ESO},pdfkeywords={monomios, polinomios, valor numérico, suma y resta, producto, identidades notables, división de polinomios, regla de Ruffini, teorema del resto, raíces, factorización, fracciones algebraicas}}

% colors de la web (bloc d'Àlgebra: blau)
\definecolor{fmnavy}{HTML}{1B2748}
\definecolor{fmblue}{HTML}{3D6AA0}
\definecolor{fmbluel}{HTML}{4E81BD}
\definecolor{fmyellow}{HTML}{FFC94D}
\definecolor{fmgold}{HTML}{E9A919}
\definecolor{fmaccent}{HTML}{FF7A50}
\definecolor{fmdef}{HTML}{1D4F73}\definecolor{fmdefbg}{HTML}{EEF4F9}
\definecolor{fmex}{HTML}{A3203A}\definecolor{fmexbg}{HTML}{FBF0F2}
\definecolor{fmsf}{HTML}{C2410C}
\definecolor{fmok}{HTML}{15803D}
\definecolor{fmmuted}{HTML}{5D6685}
\definecolor{fmsurface}{HTML}{EEF1F8}
\definecolor{fmA}{HTML}{2F6FD6}\definecolor{fmB}{HTML}{EA6A10}\definecolor{fmR}{HTML}{E11D48}\definecolor{fmK}{HTML}{16A34A}\definecolor{fmV}{HTML}{7C3AED}

% lletres
\newcommand{\comic}{\fontfamily{ComicNeue-TLF}\selectfont}
\newcommand{\heavy}{\fontseries{eb}\selectfont}
\newcommand{\semi}{\fontseries{sb}\selectfont}
\DeclareUnicodeCharacter{2212}{-}

\setlength{\parindent}{0pt}
\setlength{\parskip}{7pt}
\setlength{\lineskip}{4pt}\setlength{\lineskiplimit}{4pt}
\setlength{\emergencystretch}{3em}
\setlist{itemsep=3pt,topsep=3pt,leftmargin=1.4em}
\setcounter{secnumdepth}{1}
\renewcommand{\thesection}{3.\arabic{section}}
\titleformat{\section}[hang]{\color{fmnavy}\fontsize{23}{27}\heavy\selectfont}{\tcbox[on line,colback=fmblue,colframe=fmblue,arc=3mm,boxsep=2pt,left=6pt,right=6pt,top=3pt,bottom=3pt]{\color{white}\thesection}}{0.45em}{}[\vspace{5pt}{\color{fmblue}\titlerule[2pt]}]
\titlespacing*{\section}{0pt}{4pt}{6pt}
\titleformat{\subsection}{\color{fmblue}\Large\semi}{}{0pt}{\faCaretRight\ }
\titlespacing*{\subsection}{0pt}{8pt}{2pt}

% frase de cada apartat, «a mà»
\newcommand{\frase}[1]{\par\vspace{-2pt}{\comic\Large\color{fmbluel}\faCommentDots[regular]\ \ #1\par}\vspace{2pt}}
\newcommand{\seccio}[3]{\Needspace{16\baselineskip}\section{#2}\hypertarget{#1}{}\frase{#3}}
% notes a mà dins dels exemples
\newcommand{\nota}[1]{{\comic\small\color{fmex!85!black}#1}}
\newcommand{\notab}[1]{{\comic\small\color{fmblue}#1}}
% bé / malament
\newcommand{\be}[1]{{\color{fmok}\faCheckCircle}\ \ensuremath{#1}}
\newcommand{\mal}[1]{{\color{fmR}\faTimesCircle}\ {\color{fmR}\ensuremath{#1}}}
% pas numerat
\newcommand{\pas}[1]{\tikz[baseline=(p.base)]\node[circle,fill=fmblue,text=white,inner sep=1.2pt,minimum size=1.25em,font=\small\heavy](p){#1};}
% etiqueta de color
\newcommand{\chip}[2]{\tcbox[on line,colback=#1!14,colframe=#1!60,arc=2mm,boxsep=1pt,left=5pt,right=5pt,top=1pt,bottom=1pt]{#2}}
\newcommand{\fmkey}[1]{\colorbox{fmyellow!55}{\strut #1}}
\newcommand{\cA}[1]{{\color{fmA}#1}}\newcommand{\cB}[1]{{\color{fmB}#1}}\newcommand{\cK}[1]{{\color{fmK}#1}}\newcommand{\cR}[1]{{\color{fmR}#1}}\newcommand{\cV}[1]{{\color{fmV}#1}}

% caixes
\tcbset{fmbase/.style={enhanced,arc=3mm,boxrule=1.2pt,left=4mm,right=4mm,top=2mm,bottom=2.5mm,before skip=8pt,after skip=8pt,fonttitle=\heavy\large,before upper=\raggedright}}
\newtcolorbox{idea}[1][Idea clave]{fmbase,colback=fmdefbg,colframe=fmdef,coltitle=white,title={\faKey\ \ #1}}
\newtcolorbox{exemple}[1][Ejemplo]{fmbase,colback=fmexbg,colframe=fmex,colbacktitle=fmexbg,coltitle=fmex,title={\faPencil*\ \ #1},titlerule=0.5pt,titlerule style={fmex!40,dashed}}
\newtcolorbox{compte}[1][¡Cuidado!]{fmbase,colback=fmsf!6,colframe=fmsf,coltitle=white,fonttitle=\comic\bfseries\Large,title={\faExclamationTriangle\ \ #1}}
\newtcolorbox{truc}[1][Truco]{fmbase,colback=fmyellow!20,colframe=fmgold,colbacktitle=fmyellow,coltitle=fmnavy,fonttitle=\comic\bfseries\Large,title={\faLightbulb[regular]\ \ #1}}

% dibuix de la targeta de la unitat, en blanc
\newcommand{\fmcard}{\begin{tikzpicture}[x=0.062cm,y=-0.062cm,white,line width=3pt,line cap=round]
  \begin{scope}[transparency group,opacity=0.15]
  \node[anchor=base west,inner sep=0pt,font=\fontsize{28}{28}\selectfont] at (14,44) {$3x^2-5x+2$};
  \node[anchor=base west,inner sep=0pt,font=\fontsize{24}{24}\selectfont] at (14,100) {$(a+b)^2$};
  \end{scope}
\end{tikzpicture}}
\newcommand{\fmbrand}{\raisebox{-0.25\height}{\includegraphics[height=0.5cm]{logo-fentmates.pdf}}}

% regla de Ruffini: #1 = a, #2 = coeficients, #3 = fila del mig (sense el primer), #4 = fila de baix, #5 = nombre de coeficients
\newcommand{\ruffini}[5]{\begin{tikzpicture}[x=1.25cm,y=0.85cm,every node/.style={font=\large}]
  \foreach \v [count=\i] in {#2} \node (t\i) at (\i,0) {$\v$};
  \foreach \v [count=\i] in {#3} {\pgfmathtruncatemacro\j{\i+1}\node (m\j) at (\j,-1) {$\v$};}
  \foreach \v [count=\i] in {#4} \node (b\i) at (\i,-2.1) {$\v$};
  \node[text=fmB] at (0,-1) {$#1$};
  \draw[line width=1.1pt,fmnavy] (0.45,0.5)--(0.45,-1.5) (0.45,-1.5)--(#5+0.5,-1.5);
  \node[draw=fmR,line width=1.3pt,rounded corners=2mm,minimum size=0.75cm] at (#5,-2.1) {};
  \draw[-{Latex[length=2.2mm]},fmK,line width=1.1pt] (t1.south) -- (b1.north);
  \pgfmathtruncatemacro\last{#5-1}
  \foreach \i in {1,...,\last} {\pgfmathtruncatemacro\j{\i+1}
    \draw[-{Latex[length=2.2mm]},fmB,line width=1.1pt,shorten >=1pt] (b\i.north east) -- (m\j.south west);
    \draw[-{Latex[length=2.2mm]},fmA!70,line width=1pt,dashed,shorten >=2pt] ([xshift=8pt]t\j.south) -- ([xshift=8pt]b\j.north);}
\end{tikzpicture}}

% capçalera i peu
\pagestyle{fancy}\fancyhf{}
\fancyhead[L]{\small\fmbrand}
\fancyhead[R]{\small\color{fmmuted}4.º ESO · Unidad 3: Polinomios}
\fancyfoot[L]{\scriptsize\color{fmmuted}\href{https://creativecommons.org/licenses/by-nc-sa/4.0/deed.es}{\ccbyncsa}\ \ FentMates · fentmates.org · CC BY-NC-SA 4.0}
\fancyfoot[R]{\small\color{fmmuted}\thepage}
\renewcommand{\headrulewidth}{0.4pt}\renewcommand{\headrule}{\color{fmmuted!40}\hrule height \headrulewidth}
\fancypagestyle{primera}{\fancyhf{}\renewcommand{\headrulewidth}{0pt}
  \fancyfoot[L]{\scriptsize\color{fmmuted}\href{https://creativecommons.org/licenses/by-nc-sa/4.0/deed.es}{\ccbyncsa}\ \ FentMates · fentmates.org · CC BY-NC-SA 4.0}
  \fancyfoot[R]{\small\color{fmmuted}\thepage}}

% targeta de l'índex
\newcommand{\itile}[4]{\begin{tcolorbox}[enhanced,colback=white,colframe=fmblue!35,arc=3mm,boxrule=1pt,boxsep=0pt,left=6pt,right=6pt,top=5pt,bottom=5pt,before skip=0pt,after skip=7pt]
  \hyperlink{#1}{\begin{minipage}[c]{0.9cm}\centering{\color{fmblue}\LARGE #3}\end{minipage}\hspace{6pt}%
  \begin{minipage}[c]{0.78\linewidth}{\color{fmblue}\heavy #2}\\[-2pt]{\color{fmnavy}\semi\large #4}\end{minipage}}\end{tcolorbox}}

% taula amb vores arrodonides: el contingut es retalla amb la forma (sense @{} a les vores ni entre columnes)
\newsavebox\tabclipbox
\newcommand{\taulaclip}[3][3mm]{\sbox\tabclipbox{#3}\begin{tikzpicture}[baseline=(tcn.base)]
  \begin{scope}\clip[rounded corners=#1] (0,-\dp\tabclipbox) rectangle (\wd\tabclipbox,\ht\tabclipbox);
  \node[inner sep=0pt,outer sep=0pt,anchor=base west] (tcn) at (0,0) {\usebox\tabclipbox};\end{scope}
  \draw[#2,rounded corners=#1] (0,-\dp\tabclipbox) rectangle (\wd\tabclipbox,\ht\tabclipbox);
\end{tikzpicture}}
\newcommand{\taularod}[3][fmblue]{\taulaclip{#1,line width=1.2pt}{\renewcommand{\arraystretch}{1.5}\begin{tabular}{#2}#3\end{tabular}}}
% balança
\newcommand{\balanca}[3]{\begin{tikzpicture}[scale=#3,transform shape,every node/.style={anchor=center}]
  \fill[fmnavy] (-0.45,-1.9) rectangle (0.45,-1.75);
  \fill[fmnavy!80] (-0.12,-1.75) rectangle (0.12,0);
  \fill[fmgold] (0,0.05) circle (0.13);
  \draw[fmnavy,line width=2.4pt,line cap=round] (-2,0) -- (2,0);
  \foreach \s in {-1,1} {
    \draw[fmmuted,line width=0.8pt] (\s*2,0) -- (\s*2-0.75,-1.05) (\s*2,0) -- (\s*2+0.75,-1.05);
    \fill[fmyellow!70,draw=fmgold,line width=1pt] (\s*2-0.95,-1.05) arc[start angle=180,end angle=360,x radius=0.95,y radius=0.32] -- cycle;}
  \node[font=\Large,fill=fmA!15,draw=fmA,rounded corners=1.5mm,inner sep=3pt] at (-2,-0.68) {$#1$};
  \node[font=\Large,fill=fmB!15,draw=fmB,rounded corners=1.5mm,inner sep=3pt] at (2,-0.68) {$#2$};
\end{tikzpicture}}
% eixos amb quadrícula (dins d'un tikzpicture): xmin, xmax, ymin, ymax
\newcommand{\eixos}[4]{\draw[fmmuted!22,line width=0.4pt] (#1,#3) grid (#2,#4);
  \draw[-{Latex[length=2mm]},fmmuted,line width=0.8pt] (#1-0.2,0)--(#2+0.4,0) node[right,font=\small]{$x$};
  \draw[-{Latex[length=2mm]},fmmuted,line width=0.8pt] (0,#3-0.2)--(0,#4+0.4) node[above,font=\small]{$y$};}
% taula de traduccions (diccionari) amb les vores arrodonides
\newcommand{\dic}[2]{#1 & {\large$#2$}\\}
\newcommand{\dicratio}{0.62}
\newcommand{\dichead}{\begin{tabular}{@{}>{\raggedright}p{\dimexpr\dicratio\linewidth-12pt}>{\centering\arraybackslash}p{\dimexpr\linewidth-\dicratio\linewidth-12pt}@{}}\faComment[regular]\ Lo decimos así\dots & \dots\ así\end{tabular}}
\newenvironment{dicbox}{\begin{tcolorbox}[enhanced,arc=3mm,boxrule=1.2pt,colframe=fmblue,colback=white,coltitle=white,fonttitle=\heavy,
  title={\dichead},left=0pt,right=0pt,top=0pt,bottom=1.5mm,boxsep=0pt,lefttitle=6pt,righttitle=6pt,before skip=0pt,after skip=0pt]\rowcolors{1}{white}{fmsurface}%
  \begin{tabular}{@{\hspace{6pt}}>{\raggedright}p{\dimexpr\dicratio\linewidth-12pt}>{\centering\arraybackslash}p{\dimexpr\linewidth-\dicratio\linewidth-12pt}@{\hspace{6pt}}}}
  {\end{tabular}\end{tcolorbox}}
% resum
\newcommand{\rtile}[3]{\begin{tcolorbox}[enhanced,colback=white,colframe=fmblue!45,arc=3mm,boxrule=1pt,left=3mm,right=3mm,top=2.5mm,bottom=2.5mm,fontupper=\large,
  title={\tcbox[on line,colback=fmblue,colframe=fmblue,arc=2mm,boxsep=0pt,left=4pt,right=4pt,top=1pt,bottom=1pt]{\color{white}\small\heavy #1}\ \ #2},coltitle=fmnavy,colbacktitle=fmdefbg,fonttitle=\semi,toptitle=0.5mm,bottomtitle=0.5mm]#3\end{tcolorbox}}
% portada: #1 subtítol a mà
\newcommand{\portada}[1]{\thispagestyle{primera}
\begin{tikzpicture}[remember picture,overlay]
  \fill[fmblue] (current page.north west) rectangle ([yshift=-9.4cm]current page.north east);
  \begin{scope}
    \clip (current page.north west) rectangle ([yshift=-9.4cm]current page.north east);
    \node[anchor=north west,inner sep=0pt] at ([xshift=6.6cm,yshift=-1.3cm]current page.north west) {\fmcard};
  \end{scope}
  \node[anchor=north west,inner sep=0pt] at ([xshift=1.6cm,yshift=-1.1cm]current page.north west) {\tcbox[colback=white,colframe=white,arc=2.5mm,boxsep=0pt,left=7pt,right=7pt,top=5pt,bottom=5pt]{\includegraphics[height=0.75cm]{logo-fentmates.pdf}}};
\end{tikzpicture}
\vspace*{0.55cm}
\tcbox[colback=white!15!fmblue,colframe=white!15!fmblue,arc=2.5mm,boxsep=0pt,left=6pt,right=6pt,top=3pt,bottom=3pt]{\color{white}\small\bfseries Apuntes de teoría · 4.º ESO · Matemáticas A}

\vspace{2pt}
{\raggedright\color{white}\fontsize{30}{36}\heavy\selectfont{\color{fmyellow}Unidad 3}\\[2pt]Polinomios\par}
\vspace{8pt}
{\color{white}\comic\Large #1\par}
\vspace*{\dimexpr 7.4cm-\pagetotal\relax}
}
\newcommand{\llegenda}{\par\vspace{16pt}
{\color{fmnavy}\heavy\Large\faPalette\ \ Cómo se leen estos apuntes}\par\smallskip
\begin{tcbraster}[raster columns=4,raster equal height,raster column skip=8pt,raster before skip=12pt]
\tcbox[colback=fmdefbg,colframe=fmdef,arc=3mm,boxrule=1.2pt,halign=center,width=\linewidth]{{\color{fmdef}\Large\faKey}\\[2pt]{\color{fmdef}\heavy Idea clave}\\[-1pt]{\small lo esencial}}
\tcbox[colback=fmexbg,colframe=fmex,arc=3mm,boxrule=1.2pt,halign=center,width=\linewidth]{{\color{fmex}\Large\faPencil*}\\[2pt]{\color{fmex}\heavy Ejemplo}\\[-1pt]{\small paso a paso}}
\tcbox[colback=fmsf!6,colframe=fmsf,arc=3mm,boxrule=1.2pt,halign=center,width=\linewidth]{{\color{fmsf}\Large\faExclamationTriangle}\\[2pt]{\color{fmsf}\heavy ¡Cuidado!}\\[-1pt]{\small errores típicos}}
\tcbox[colback=fmyellow!20,colframe=fmgold,arc=3mm,boxrule=1.2pt,halign=center,width=\linewidth]{{\color{fmgold!80!black}\Large\faLightbulb[regular]}\\[2pt]{\color{fmnavy}\heavy Truco}\\[-1pt]{\small para recordarlo}}
\end{tcbraster}}
\newcommand{\resumtitol}{\clearpage
\hypertarget{resum}{}
{\color{fmnavy}\fontsize{23}{27}\heavy\selectfont\faClipboardList\ \ En resumen}\par\vspace{5pt}{\color{fmblue}\titlerule[2pt]}
\frase{Toda la unidad en una página.}}
% els 5 passos per a resoldre problemes: #2 = text del pas 2
\newcommand{\passosproblema}[2]{\begin{center}
\begin{tikzpicture}[s/.style={draw=fmblue,fill=fmdefbg,rounded corners=3mm,line width=1.2pt,minimum width=2.75cm,minimum height=2.45cm,align=center,font=\small},
  ic/.style={font=\LARGE,text=fmblue}]
  \foreach \i/\ic/\t in {1/\faBookOpen/{\textbf{Leemos}\\y entendemos},2/\faQuestionCircle/{#1},3/\faEquals/{#2},4/\faCogs/{\textbf{Resolvemos}},5/\faCheckDouble/{\textbf{Comprobamos}\\y respondemos}} {
    \node[s] (s\i) at (3.3*\i,0) {};
    \node[ic] at ([yshift=-0.5cm]s\i.north) {\ic};
    \node[align=center,font=\small,anchor=north] at ([yshift=-0.95cm]s\i.north) {\t};
    \node[circle,fill=fmblue,text=white,font=\small\heavy,inner sep=1.5pt,minimum size=1.3em] at (s\i.north west) {\i};}
  \foreach \i [evaluate=\i as \j using int(\i+1)] in {1,...,4} \draw[-{Latex[length=3mm]},fmmuted,line width=1.5pt] (s\i) -- (s\j);
\end{tikzpicture}
\end{center}}

\newcommand{\lo}{\mathopen{]}}\newcommand{\ro}{\mathclose{[}}
% punts ple i buit, nombres dels eixos i taula de valors
\newcommand{\pple}[2][fmblue]{\fill[#1] #2 circle (3.6pt);}
\newcommand{\nums}[4]{\foreach \k in {#1,...,#2} {\ifnum\k=0\else\node[below,font=\scriptsize,text=fmmuted] at (\k,0) {$\k$};\fi}
  \foreach \k in {#3,...,#4} {\ifnum\k=0\else\node[left,font=\scriptsize,text=fmmuted] at (0,\k) {$\k$};\fi}}
\newcommand{\tv}[3][fmblue]{\taulaclip{#1,line width=1.2pt}{\renewcommand{\arraystretch}{1.45}%
  \begin{tabular}{>{\columncolor{#1}\color{white}\heavy}c*{12}{w{c}{0.8cm}}}#2\\\rowcolor{#1!8}\cellcolor{#1}#3\end{tabular}}}
% etiqueta d'una recta o corba
\newcommand{\etq}[3]{\node[fill=white,inner sep=1.5pt,text=#1,font=\small] at #2 {#3};}

% interval sobre la recta: #1 esquerre, #2 dret, #3 tipus esquerre (o/c/a), #4 tipus dret (o/c/a), #5 color
\def\intvdot#1#2#3{\ifx#2o\filldraw[fill=white,draw=#3,line width=1.3pt] (#1,0) circle (3.3pt);\fi\ifx#2c\fill[#3] (#1,0) circle (3.5pt);\fi}
\newcommand{\intv}[5]{\begin{tikzpicture}[x=0.4cm,baseline=-3pt]
  \draw[fmmuted,line width=0.8pt] (-6,0)--(7,0);
  \foreach \k in {-5,...,6} \draw[fmmuted,line width=0.6pt] (\k,-0.08)--(\k,0.08);
  \ifx#3a\def\xa{-6}\else\def\xa{#1}\fi
  \ifx#4a\def\xb{7}\else\def\xb{#2}\fi
  \ifx#4a\draw[#5,line width=2.6pt,-{Latex[length=2.6mm]}] (\xa,0)--(\xb,0);\else\ifx#3a\draw[#5,line width=2.6pt,{Latex[length=2.6mm]}-] (\xa,0)--(\xb,0);\else\draw[#5,line width=2.6pt] (\xa,0)--(\xb,0);\fi\fi
  \intvdot{#1}{#3}{#5}\intvdot{#2}{#4}{#5}
  \ifx#3a\else\node[below=2pt,font=\scriptsize] at (#1,0) {$#1$};\fi
  \ifx#4a\else\node[below=2pt,font=\scriptsize] at (#2,0) {$#2$};\fi
\end{tikzpicture}}

\newcommand{\tsq}[2]{\filldraw[fill=fmA!30,draw=fmA,line width=0.8pt,rounded corners=0.5pt] (#1,#2) rectangle ++(0.9,0.9); \node[font=\small] at ($(#1,#2)+(0.45,0.45)$) {$x^2$};}


\begin{document}
\portada{Expresiones con letras: operamos, dividimos y factorizamos.}

{\color{fmnavy}\heavy\Large\faMapSigns\ \ ¿Qué aprenderás?}\par\smallskip
\begin{minipage}[t]{0.485\linewidth}
\itile{sec1}{3.1}{\faFont}{Monomios}
\itile{sec2}{3.2}{\faEquals}{Polinomios. Valor numérico}
\itile{sec3}{3.3}{\faPlus}{Suma y resta}
\itile{sec4}{3.4}{\faTimes}{Producto. Identidades notables}
\itile{sec5}{3.5}{\faDivide}{División de polinomios}
\end{minipage}\hfill
\begin{minipage}[t]{0.485\linewidth}
\itile{sec6}{3.6}{\faTable}{Regla de Ruffini}
\itile{sec7}{3.7}{\faSearch}{Raíces de un polinomio}
\itile{sec8}{3.8}{\faPuzzlePiece}{Factorización}
\itile{sec9}{3.9}{\faIcon{percentage}}{Fracciones algebraicas}
\itile{resum}{\phantom{3.9}}{\faClipboardList}{En resumen}
\end{minipage}
\llegenda

% ============================================================ 3.1
\clearpage
\seccio{sec1}{Monomios}{Un monomio es un número multiplicado por letras.}

\vfill
\begin{center}
\begin{tikzpicture}
  \node[font=\fontsize{48}{48}\selectfont,inner sep=2pt] (c) {$\cA{-5}$};
  \node[font=\fontsize{48}{48}\selectfont,inner sep=2pt,right=2pt of c] (l) {$\cB{x^{\cK{2}}y^{\cK{1}}}$};
  \node[below left=0.9cm and -0.3cm of c,fill=fmA!14,draw=fmA,rounded corners=2mm,text=fmA,font=\heavy\large] (cl) {coeficiente};
  \node[below right=0.9cm and -0.6cm of l,fill=fmB!14,draw=fmB,rounded corners=2mm,text=fmB,font=\heavy\large] (ll) {parte literal};
  \node[right=2.4cm of l,fill=fmK!14,draw=fmK,rounded corners=2mm,text=fmK!80!black,font=\large,align=center] (g) {{\heavy grado}\\ $2+1=\mathbf{3}$};
  \draw[-{Latex[length=3mm]},fmA,line width=1.3pt] (cl) -- (c);
  \draw[-{Latex[length=3mm]},fmB,line width=1.3pt] (ll) -- (l);
  \draw[-{Latex[length=3mm]},fmK,line width=1.3pt] (g.west) to[bend right=25] ([yshift=-2pt]l.north east);
  \node[font=\comic,text=fmmuted,above=0pt of g] {suma los exponentes};
  \node[font=\comic\small,text=fmmuted,below=2pt of g] {si no hay exponente, es 1: \ $y=y^1$};
\end{tikzpicture}
\end{center}

\vfill
\begin{minipage}[c]{0.43\linewidth}
\begin{idea}
\begin{itemize}[leftmargin=1em]
\item \cA{\textbf{Coeficiente}}: el número (con el signo).
\item \cB{\textbf{Parte literal}}: las letras y los exponentes.
\item \cK{\textbf{Grado}}: la suma de los exponentes.
\end{itemize}
\end{idea}
\end{minipage}\hfill
\begin{minipage}[c]{0.54\linewidth}\centering\renewcommand{\arraystretch}{1.35}
\taularod{*{4}{>{\hspace{2pt}}c<{\hspace{2pt}}}}{
\rowcolor{fmblue}\color{white}\heavy Monomio & \color{white}\heavy Coeficiente & \color{white}\heavy Parte literal & \color{white}\heavy Grado\\
$7x^4$ & \cA{$7$} & \cB{$x^4$} & \cK{4}\\
\rowcolor{fmsurface}$-a^2b^3$ & \cA{$-1$} & \cB{$a^2b^3$} & \cK{$2+3=5$}\\
$x$ & \cA{$1$} & \cB{$x$} & \cK{1}\\
\rowcolor{fmsurface}$-6$ & \cA{$-6$} & \cB{---} & \cK{0}
}
\end{minipage}

\vfill
\begin{minipage}[t]{0.44\linewidth}
\begin{truc}
¿No ves el coeficiente? Es un \textbf{1}:\\
$x^2=\cA{1}x^2$ \qquad $-x=\cA{-1}x$\\[2pt]
Un número solo es un monomio de \textbf{grado 0}.
\end{truc}
\end{minipage}\hfill
\begin{minipage}[t]{0.53\linewidth}
\begin{compte}[Esto NO son monomios]
\mal{\dfrac{3}{x}}\ \ {\small la letra divide}\\[3pt]
\mal{\sqrt{x}}\ \ {\small hay una raíz}\\[3pt]
\mal{2x+3}\ \ {\small es una suma (un polinomio)}
\end{compte}
\end{minipage}

\vfill
\subsection{Monomios semejantes: la misma familia}
{\large Son \textbf{semejantes} si tienen \textbf{la misma parte literal}} (mismas letras y mismos exponentes). Es como tener el mismo apellido.

\vfill
\begin{center}
\begin{tikzpicture}[card/.style={draw=#1,fill=#1!10,line width=1.2pt,rounded corners=3mm,minimum width=2.6cm,minimum height=1.25cm,font=\Large}]
  \node[card=fmK] (a) at (0,0) {$3\cB{x^2y}$};
  \node[card=fmK] (b) at (3.4,0) {$-\frac12\cB{x^2y}$};
  \node[font=\Large,text=fmK] at (1.7,0) {\faUsers};
  \node[text=fmK!80!black,font=\comic\large] at (1.7,-1.05) {\faCheckCircle\ sí son semejantes};
  \node[card=fmR] (c) at (8.4,0) {$3\cB{x^2y}$};
  \node[card=fmR] (d) at (11.8,0) {$3\cB{xy^2}$};
  \node[font=\Large,text=fmR] at (10.1,0) {$\neq$};
  \node[text=fmR,font=\comic\large] at (10.1,-1.05) {\faTimesCircle\ no lo son (exponentes cambiados)};
\end{tikzpicture}
\end{center}
\begin{idea}[¿Para qué sirve?]
\centering\Large Solo se pueden \textbf{sumar y restar} los monomios \textbf{semejantes}.
\end{idea}

\vfill
\clearpage

\vfill
\subsection{Suma y resta}
\begin{minipage}[c]{0.48\linewidth}
\begin{tikzpicture}[font=\Large]
  \node at (0,0) {3 \textcolor{fmR}{\faIcon{apple-alt}}};
  \node at (1.15,0) {$+$};
  \node at (2.3,0) {2 \textcolor{fmR}{\faIcon{apple-alt}}};
  \node at (3.45,0) {$=$};
  \node at (4.6,0) {\textbf{5} \textcolor{fmR}{\faIcon{apple-alt}}};
  \node[text=fmok] at (5.7,0) {\faCheckCircle};
  \node at (0,-1.1) {3 \textcolor{fmR}{\faIcon{apple-alt}}};
  \node at (1.15,-1.1) {$+$};
  \node at (2.3,-1.1) {2 \textcolor{fmK}{\faLemon}};
  \node at (3.45,-1.1) {$\to$};
  \node[font=\comic\normalsize,text=fmR,align=center] at (4.75,-1.1) {no se pueden\\juntar};
  \node[text=fmR] at (5.95,-1.1) {\faTimesCircle};
\end{tikzpicture}
\end{minipage}\hfill
\begin{minipage}[c]{0.5\linewidth}
\begin{tikzpicture}[x=1cm,y=1cm]
  \foreach \i in {0,...,4} {\tsq{\i*1}{0}}
  \node[font=\Large] at (5.25,0.45) {$+$};
  \foreach \i in {0,...,2} {\tsq{5.6+\i*1}{0}}
  \node[font=\Large,anchor=west] at (0,-0.65) {$5x^2+3x^2=\mathbf{8x^2}$};
  \node[font=\comic,text=fmmuted,anchor=west] at (4.6,-0.65) {8 piezas iguales};
\end{tikzpicture}
\end{minipage}

\vfill
\begin{minipage}[t]{0.52\linewidth}
\begin{idea}
\textbf{Sumamos los coeficientes}. La parte literal se queda \textbf{igual}.\\[4pt]
$5x^2+3x^2=(5+3)x^2=8x^2$\\[2pt]
$4ab-7ab=-3ab$\\[2pt]
$5x^2+3x$ \ \notab{no son semejantes: se queda así}
\end{idea}
\end{minipage}\hfill
\begin{minipage}[t]{0.45\linewidth}
\begin{compte}[Error típico]
\mal{5x^2+3x^2=8x^4}\\[3pt]
\be{5x^2+3x^2=8x^2}\\[3pt]
Cuando \textbf{sumamos}, los exponentes \textbf{no cambian}.
\end{compte}
\end{minipage}

\vfill
\subsection{Producto, cociente y potencia}
\begin{tcbraster}[raster columns=3,raster equal height,raster column skip=8pt]
\begin{tcolorbox}[fmbase,colback=fmdefbg,colframe=fmdef,coltitle=white,title={\faTimes\ \ Producto},before upper=\centering]
Números: \textbf{multiplicamos}\\ Exponentes: \fmkey{\textbf{sumamos}}\\[3pt]
{\large$x^m\cdot x^n=x^{m+n}$}\\[4pt]
$(-3x^2y)\cdot(4xy^3)=-12x^3y^4$
\end{tcolorbox}
\begin{tcolorbox}[fmbase,colback=fmdefbg,colframe=fmdef,coltitle=white,title={\faDivide\ \ Cociente},before upper=\centering]
Números: \textbf{dividimos}\\ Exponentes: \fmkey{\textbf{restamos}}\\[3pt]
{\large$x^m:x^n=x^{m-n}$}\\[4pt]
$15x^5y^2:(3x^2y)=5x^3y$
\end{tcolorbox}
\begin{tcolorbox}[fmbase,colback=fmdefbg,colframe=fmdef,coltitle=white,title={\faSuperscript\ \ Potencia},before upper=\centering]
Número: \textbf{lo elevamos}\\ Exponentes: \fmkey{\textbf{multiplicamos}}\\[3pt]
{\large$(x^m)^n=x^{m\cdot n}$}\\[4pt]
$(2x^3)^2=4x^6$
\end{tcolorbox}
\end{tcbraster}

\vfill
\begin{minipage}[t]{0.62\linewidth}
\begin{exemple}
\vspace{-6pt}
\begin{align*}
2x^3\cdot5x^4&=(\cA{2\cdot5})\,x^{\cK{3+4}}=\boxed{10x^7}\\
20x^6y^3:(-4x^2y)&=(\cA{20:(-4)})\,x^{\cK{6-2}}y^{\cK{3-1}}=\boxed{-5x^4y^2}
\end{align*}
\nota{\cA{números con números} \quad\cK{exponentes de cada letra}}
\end{exemple}
\end{minipage}\hfill
\begin{minipage}[t]{0.35\linewidth}
\begin{truc}
Recuerda: \ $x^0=1$\\[4pt]
$6x^2y:(2x^2)=3\,\cR{x^0}\,y$\\[2pt]
$\phantom{6x^2y:(2x^2)}=3\cdot\cR{1}\cdot y=3y$
\end{truc}
\end{minipage}

\vfill
% ============================================================ 3.2
\clearpage
\seccio{sec2}{Polinomios. Valor numérico}{Un polinomio es una suma de monomios (no semejantes).}

\vfill
\begin{center}
\chip{fmblue}{\semi Monomio} \ 1 término: $5x$ \qquad
\chip{fmblue}{\semi Binomio} \ 2 términos: $2x+3$ \qquad
\chip{fmblue}{\semi Trinomio} \ 3 términos: $x^2-5x+6$
\end{center}

\vfill
\begin{center}
\begin{tikzpicture}[t/.style={font=\fontsize{30}{30}\selectfont,inner sep=3pt,text height=25pt,text depth=8pt}]
  \node[t] (P) {$P(x)=$};
  \node[t,right=0pt of P,fill=fmA!12,rounded corners=2mm] (a) {$\cA{4}x^{\cK{3}}$};
  \node[t,right=5pt of a] (b) {$-2x^2$};
  \node[t,right=2pt of b] (c) {$+x$};
  \node[t,right=2pt of c,fill=fmB!14,rounded corners=2mm] (d) {$\cB{-7}$};
  \node[above left=0.7cm and -0.6cm of a,fill=fmA!14,draw=fmA,rounded corners=2mm,font=\large,align=center] (cp) {{\heavy\color{fmA}coeficiente principal}: 4};
  \coordinate (e3) at ([xshift=-10pt,yshift=-19pt]a.north east);% sota l'exponent 3
  \node[anchor=north,fill=fmK!14,draw=fmK,rounded corners=2mm,font=\large,align=center] (gr) at ([yshift=-0.8cm]e3 |- a.south) {{\heavy\color{fmK!80!black}grado del polinomio}: 3\\[-2pt]{\comic\small el grado más alto}};
  \node[above right=0.7cm and -0.8cm of d,fill=fmB!14,draw=fmB,rounded corners=2mm,font=\large,align=center] (ti) {{\heavy\color{fmB}término independiente}: $-7$\\[-2pt]{\comic\small el que no tiene letra}};
  \draw[-{Latex[length=3mm]},fmA,line width=1.3pt] (cp.south -| {$(a.west)+(11pt,0)$}) -- ([xshift=11pt,yshift=-1pt]a.north west);
  \draw[-{Latex[length=3mm]},fmK,line width=1.3pt] (gr.north) -- (e3);
  \draw[-{Latex[length=3mm]},fmB,line width=1.3pt] (ti) -- (d.north);
\end{tikzpicture}
\end{center}

\vfill
\begin{minipage}[t]{0.38\linewidth}
\begin{idea}[Ordenado y completo]
\textbf{Ordenado}: de mayor a menor grado.\\[3pt]
\textbf{Completo}: no falta ningún grado (hasta el término independiente).
\end{idea}
\end{minipage}\hfill
\begin{minipage}[t]{0.59\linewidth}
\begin{exemple}[Ejemplo: \ $x^4-3x+1$]
\centering
\begin{tikzpicture}[s/.style={font=\Large,minimum height=0.9cm,minimum width=1.3cm,rounded corners=2mm}]
  \node[s,draw=fmblue,fill=fmdefbg] at (0,0) {$x^4$};
  \node[s,draw=fmR,dashed,line width=1pt,text=fmR] at (1.55,0) {?\,$x^3$};
  \node[s,draw=fmR,dashed,line width=1pt,text=fmR] at (3.1,0) {?\,$x^2$};
  \node[s,draw=fmblue,fill=fmdefbg] at (4.65,0) {$-3x$};
  \node[s,draw=fmblue,fill=fmdefbg] at (6.2,0) {$+1$};
\end{tikzpicture}\\[3pt]
\raggedright Está ordenado \be{}, pero \textbf{no es completo}: faltan $x^3$ y $x^2$. \ \nota{Lo podemos escribir: $x^4+0x^3+0x^2-3x+1$}
\end{exemple}
\end{minipage}

\vfill
\subsection{Valor numérico y raíces}
{\large $P(a)$ quiere decir: cambia la $x$ por $a$ y calcula.}

\vfill
\begin{minipage}[c]{0.45\linewidth}
\begin{idea}
Si \ $P(a)=0$, \ decimos que $a$ es una \fmkey{\textbf{raíz}} del polinomio.
\end{idea}
\end{minipage}\hfill
\begin{minipage}[c]{0.52\linewidth}
\begin{exemple}[Ejemplo: \ $P(x)=x^2-5x+6$]
$P(\cR{1})=1-5+6=2$\\[3pt]
$P(\cR{2})=4-10+6=\boxed{0}$ \ \ {\color{fmok}\Large\faBullseye}\ \ \nota{¡2 es una raíz!}
\end{exemple}
\end{minipage}

\vfill
% ============================================================ 3.3
\clearpage
\seccio{sec3}{Suma y resta}{Juntamos los términos de la misma familia.}

\vfill
\begin{idea}
\begin{tabular}{@{}p{0.08\linewidth}p{0.84\linewidth}@{}}
{\color{fmblue}\Huge\faPlus} & \textbf{Suma}: ponemos un polinomio al lado del otro y \textbf{reducimos los semejantes}.\\[6pt]
{\color{fmblue}\Huge\faMinus} & \textbf{Resta}: \textbf{cambiamos el signo de TODOS} los términos del segundo y sumamos.
\end{tabular}
\end{idea}

\vfill
{\large Lo hacemos en \textbf{columnas}: cada columna, una familia.} \quad $P(x)=3x^3-2x^2+5$ \quad $Q(x)=x^3+4x^2-3x-1$

\vfill
\begin{center}\renewcommand{\arraystretch}{1.5}\large
\begin{tabular}{r>{\columncolor{fmV!10}}r>{\columncolor{fmA!10}}r>{\columncolor{fmK!10}}r>{\columncolor{fmB!10}}r}
 & {\comic\normalsize\color{fmV} las $x^3$} & {\comic\normalsize\color{fmA} las $x^2$} & {\comic\normalsize\color{fmK} las $x$} & {\comic\normalsize\color{fmB} los números}\\
$P(x)$: & $3x^3$ & $-2x^2$ & & $+5$\\
$Q(x)$: & $x^3$ & $+4x^2$ & $-3x$ & $-1$\\ \hline
$P+Q$: & $\mathbf{4x^3}$ & $\mathbf{+2x^2}$ & $\mathbf{-3x}$ & $\mathbf{+4}$
\end{tabular}
\end{center}

\vfill
\subsection{Para restar: el signo menos «entra» en el paréntesis}
\begin{center}
\begin{tikzpicture}[t/.style={font=\fontsize{22}{22}\selectfont,inner sep=2pt,text height=18pt,text depth=6pt}]
  \node[t,text=fmR] (m) {$-$};
  \node[t,right=0pt of m] (p1) {$\big($};
  \node[t,right=0pt of p1,minimum width=1.65cm] (t1) {$x^3$};
  \node[t,right=0pt of t1,minimum width=1.65cm] (t2) {$+4x^2$};
  \node[t,right=0pt of t2,minimum width=1.65cm] (t3) {$-3x$};
  \node[t,right=0pt of t3,minimum width=1.65cm] (t4) {$-1$};
  \node[t,right=0pt of t4] (p2) {$\big)$};
  \foreach \k/\b in {t1/50,t2/45,t3/40,t4/35} \draw[-{Latex[length=2.5mm]},fmR,line width=1.1pt] (m.north) to[bend left=\b] (\k.north);
  \node[t,below=1.0cm of t1,text=fmR] (r1) {$-x^3$};
  \node[t,below=1.0cm of t2,text=fmR] (r2) {$-4x^2$};
  \node[t,below=1.0cm of t3,text=fmR] (r3) {$+3x$};
  \node[t,below=1.0cm of t4,text=fmR] (r4) {$+1$};
  \foreach \i in {1,...,4} \draw[-{Latex[length=2.5mm]},fmmuted,line width=1pt] (t\i) -- (r\i);
  \node[t] at (m |- r1) {$=$};
  \node[right=1.2cm of p2,font=\comic\large,text=fmR,align=left] {cambian\\TODOS los signos};
\end{tikzpicture}
\end{center}

\vfill
\begin{minipage}[t]{0.58\linewidth}
\begin{exemple}[Ejemplo: \ $P-Q$]
\vspace{-6pt}
\begin{align*}
P-Q&=3x^3-2x^2+5\,\cR{-\,x^3-4x^2+3x+1} && \text{\nota{cambiamos signos}}\\
&=\boxed{2x^3-6x^2+3x+6} && \text{\nota{reducimos}}
\end{align*}
\end{exemple}
\end{minipage}\hfill
\begin{minipage}[t]{0.39\linewidth}
\begin{compte}[El error más frecuente]
Cambiar solo el \textbf{primer} signo:\\[3pt]
\mal{-(x^3+4x^2\dots)=-x^3+4x^2\dots}\\[3pt]
El menos afecta a \textbf{todo} el paréntesis.
\end{compte}
\end{minipage}

\vfill
% ============================================================ 3.4
\clearpage
\seccio{sec4}{Producto. Identidades notables}{Cada término multiplica TODOS los términos del otro.}

\vfill
\subsection{Monomio por polinomio}
\begin{center}
\begin{tikzpicture}[t/.style={font=\fontsize{22}{22}\selectfont,inner sep=2pt,text height=18pt,text depth=6pt}]
  \node[t,text=fmB] (m) {$-2x$};
  \node[t,right=0pt of m] (p1) {$\big($};
  \node[t,right=0pt of p1] (t1) {$3x^2$};
  \node[t,right=4pt of t1] (t2) {$-x$};
  \node[t,right=4pt of t2] (t3) {$+4$};
  \node[t,right=0pt of t3] (p2) {$\big)$};
  \foreach \k/\b/\c in {t1/55/fmV,t2/50/fmA,t3/45/fmK} \draw[-{Latex[length=2.5mm]},\c,line width=1.3pt] (m.north) to[bend left=\b] (\k.north);
  \node[t,right=0.3cm of p2] (eq) {$=$};
  \node[t,right=0.3cm of eq] {$\cV{-6x^3}\cA{+2x^2}\cK{-8x}$};
\end{tikzpicture}
\end{center}

\vfill
\subsection{Polinomio por polinomio: con una tabla}
\begin{minipage}[c]{0.5\linewidth}\centering
\begin{tikzpicture}[x=1.75cm,y=0.95cm,c/.style={font=\large}]
  \begin{scope}
  \clip[rounded corners=3mm] (0,0) rectangle (4,3);
  \fill[fmdefbg] (0,0) rectangle (1,3) (1,2) rectangle (4,3);
  \fill[fmV!18] (1,1) rectangle (2,2);
  \fill[fmA!18] (2,1) rectangle (3,2);
  \fill[fmK!18] (3,1) rectangle (4,2);
  \fill[fmA!18] (1,0) rectangle (2,1);
  \fill[fmK!18] (2,0) rectangle (3,1);
  \fill[fmB!18] (3,0) rectangle (4,1);
  \end{scope}
  \foreach \x in {1,2,3} \draw[fmblue,line width=0.8pt] (\x,0)--(\x,3);
  \foreach \y in {1,2} \draw[fmblue,line width=0.8pt] (0,\y)--(4,\y);
  \draw[fmblue,line width=1.2pt,rounded corners=3mm] (0,0) rectangle (4,3);
  \node[c,font=\large\heavy,text=fmblue] at (0.5,2.5) {$\times$};
  \foreach \h [count=\i] in {x^2,+4x,-1} \node[c,font=\large\heavy,text=fmB] at (\i+0.5,2.5) {$\h$};
  \foreach \h [count=\j] in {2x,-3} \node[c,font=\large\heavy,text=fmB] at (0.5,2.5-\j) {$\h$};
  \node[c] at (1.5,1.5) {$2x^3$};  \node[c] at (2.5,1.5) {$8x^2$};   \node[c] at (3.5,1.5) {$-2x$};
  \node[c] at (1.5,0.5) {$-3x^2$}; \node[c] at (2.5,0.5) {$-12x$};  \node[c] at (3.5,0.5) {$+3$};
\end{tikzpicture}\\[2pt]
{\comic\color{fmmuted}$(2x-3)\cdot(x^2+4x-1)$}
\end{minipage}\hfill
\begin{minipage}[c]{0.47\linewidth}
\pas{1} Cada casilla: \textbf{fila $\times$ columna}.\\[4pt]
\pas{2} Juntamos las casillas del \textbf{mismo color} (son semejantes):\\[4pt]
\hspace*{1.5em}$\cV{2x^3}\ \cA{+8x^2-3x^2}\ \cK{-2x-12x}\ \cB{+3}$\\[6pt]
\pas{3} Resultado:\\[2pt]
\hspace*{1.5em}{\Large$\boxed{2x^3+5x^2-14x+3}$}
\end{minipage}

\medskip
\begin{minipage}[t]{0.5\linewidth}
\begin{exemple}[Ejemplo: \ $(x+4)(x-6)$]
\vspace{-6pt}
\begin{align*}
(x+4)(x-6)&=x^2-6x+4x-24 &\ & \text{\nota{todos por todos}}\\
&=\boxed{x^2-2x-24} &\ & \text{\nota{reducimos}}
\end{align*}
\end{exemple}
\end{minipage}\hfill
\begin{minipage}[t]{0.47\linewidth}
\begin{truc}
El \textbf{grado del producto} es la \textbf{suma de los grados}:\\[3pt]
grado 1 \ $\times$ \ grado 2 \ $\to$ \ grado \textbf{3}
\end{truc}
\end{minipage}

\vfill
\subsection{Polinomio por polinomio: término a término}
\begin{center}
\begin{tikzpicture}[t/.style={font=\fontsize{22}{22}\selectfont,inner sep=2pt,text height=18pt,text depth=6pt}]
  \node[t] (q1) {$\big($};
  \node[t,right=0pt of q1,text=fmA] (f1) {$2x$};
  \node[t,right=4pt of f1,text=fmB] (f2) {$+3$};
  \node[t,right=0pt of f2] (q2) {$\big)$};
  \node[t,right=0pt of q2] (p1) {$\big($};
  \node[t,right=0pt of p1] (t1) {$x^2$};
  \node[t,right=4pt of t1] (t2) {$-x$};
  \node[t,right=4pt of t2] (t3) {$+4$};
  \node[t,right=0pt of t3] (p2) {$\big)$};
  \foreach \k/\b in {t1/50,t2/45,t3/40} \draw[-{Latex[length=2.5mm]},fmA,line width=1.3pt] (f1.north) to[bend left=\b] (\k.north);
  \foreach \k/\b in {t1/50,t2/45,t3/40} \draw[-{Latex[length=2.5mm]},fmB,line width=1.3pt] (f2.south) to[bend right=\b] (\k.south);
  \node[font=\comic,text=fmA] at ([yshift=1.35cm]t2.north) {el $2x$ por todos};
  \node[font=\comic,text=fmB] at ([yshift=-1.35cm]t2.south) {el $+3$ por todos};
  \node[anchor=west,font=\Large] (r1) at ([xshift=0.5cm,yshift=0.5cm]p2.east) {$=\cA{2x^3-2x^2+8x}\,\cB{+3x^2-3x+12}$};
  \node[anchor=north west,font=\Large,inner sep=0pt] (r2) at ([yshift=-8pt]r1.south west) {\hspace{3.333pt}$=\boxed{2x^3+x^2+5x+12}$};
  \node[anchor=west,font=\comic\small,text=fmex!85!black] at (r2.east) {\ \ reducimos los semejantes};
\end{tikzpicture}
\end{center}

\vfill
\clearpage
\subsection{Identidades notables}

\vfill
\begin{tcbraster}[raster columns=3,raster equal height,raster column skip=8pt]
\begin{tcolorbox}[fmbase,colback=fmdefbg,colframe=fmdef,coltitle=white,title={\faStar\ \ Cuadrado de una suma},before upper=\centering,fonttitle=\heavy\small,left=2mm,right=2mm]
\vspace{4pt}{\large$(\cA{a}+\cB{b})^2=\cA{a}^2\,{+}\,2\cA{a}\cB{b}+\cB{b}^2$}
\end{tcolorbox}
\begin{tcolorbox}[fmbase,colback=fmdefbg,colframe=fmdef,coltitle=white,title={\faStar\ \ Cuadrado de una diferencia},before upper=\centering,fonttitle=\heavy\small,left=2mm,right=2mm]
\vspace{4pt}{\large$(\cA{a}-\cB{b})^2=\cA{a}^2\,{-}\,2\cA{a}\cB{b}+\cB{b}^2$}
\end{tcolorbox}
\begin{tcolorbox}[fmbase,colback=fmdefbg,colframe=fmdef,coltitle=white,title={\faStar\ \ Suma por diferencia},before upper=\centering,fonttitle=\heavy\small,left=2mm,right=2mm]
\vspace{4pt}{\large$(\cA{a}+\cB{b})(\cA{a}-\cB{b})=\cA{a}^2-\cB{b}^2$}
\end{tcolorbox}
\end{tcbraster}
\begin{center}
{\comic\large\color{fmblue}«El cuadrado del primero, más el doble del primero por el segundo, más el cuadrado del segundo»}
\end{center}

\vfill
\subsection{Lo vemos con dibujos}
\begin{minipage}[c]{0.36\linewidth}\centering
\begin{tikzpicture}[x=1cm,y=1cm]
  \def\a{2.5}\def\b{1.1}
  \fill[fmA!28] (0,\b) rectangle (\a,\a+\b); \node[font=\Large] at (\a/2,\b+\a/2) {$\cA{a}^2$};
  \fill[fmK!25] (\a,\b) rectangle (\a+\b,\a+\b); \node at (\a+\b/2,\b+\a/2) {$\cA{a}\cB{b}$};
  \fill[fmK!25] (0,0) rectangle (\a,\b); \node at (\a/2,\b/2) {$\cA{a}\cB{b}$};
  \fill[fmB!28] (\a,0) rectangle (\a+\b,\b); \node at (\a+\b/2,\b/2) {$\cB{b}^2$};
  \draw[fmnavy,line width=1pt] (0,0) rectangle (\a+\b,\a+\b) (\a,0)--(\a,\a+\b) (0,\b)--(\a+\b,\b);
  \draw[decorate,decoration={brace,amplitude=5pt},fmA,thick] (0,\a+\b+0.08)--(\a,\a+\b+0.08) node[midway,above=5pt,text=fmA]{$a$};
  \draw[decorate,decoration={brace,amplitude=5pt},fmB,thick] (\a,\a+\b+0.08)--(\a+\b,\a+\b+0.08) node[midway,above=5pt,text=fmB]{$b$};
\end{tikzpicture}\\[3pt]
{\comic\color{fmmuted}$(a+b)^2=a^2+\mathbf{2}ab+b^2$}
\end{minipage}\hfill
\begin{minipage}[c]{0.62\linewidth}\centering
\begin{tikzpicture}[x=1cm,y=1cm]
  \def\a{2.5}\def\b{1.0}
  \fill[fmA!28] (0,0) rectangle (\a,\a-\b);
  \fill[fmK!25] (0,\a-\b) rectangle (\a-\b,\a);
  \draw[fmB,dashed,line width=1pt] (\a-\b,\a-\b) rectangle (\a,\a); \node[text=fmB] at (\a-\b/2,\a-\b/2) {$b^2$};
  \draw[fmnavy,line width=1pt] (0,0)--(\a,0)--(\a,\a-\b)--(\a-\b,\a-\b)--(\a-\b,\a)--(0,\a)--cycle;
  \node[font=\comic,text=fmmuted] at (\a/2,-0.4) {$a^2-b^2$};
  \draw[-{Latex[length=3mm]},fmmuted,line width=1.4pt] (\a+0.4,\a/2) -- node[above,font=\comic\small,text=fmmuted,align=center]{movemos\\la pieza verde} (\a+1.9,\a/2);
  \begin{scope}[xshift=4.8cm]
  \fill[fmA!28] (0,0) rectangle (\a,\a-\b);
  \fill[fmK!25] (\a,0) rectangle (\a+\b,\a-\b);
  \draw[fmnavy,line width=1pt] (0,0) rectangle (\a+\b,\a-\b) (\a,0)--(\a,\a-\b);
  \draw[decorate,decoration={brace,amplitude=5pt},thick] (0,\a-\b+0.08)--(\a+\b,\a-\b+0.08) node[midway,above=5pt]{$a+b$};
  \draw[decorate,decoration={brace,amplitude=5pt,mirror},thick] (\a+\b+0.08,0)--(\a+\b+0.08,\a-\b) node[midway,right=5pt]{$a-b$};
  \node[font=\comic,text=fmmuted] at (\a/2+\b/2,-0.4) {$(a+b)(a-b)$};
  \end{scope}
\end{tikzpicture}
\end{minipage}

\vfill
\subsection{Cómo se aplican}
\begin{exemple}[Ejemplo: \ primero buscamos quién hace de \cA{$a$} y quién de \cB{$b$}]
\vspace{-6pt}
\begin{align*}
(\cA{3x}+\cB{2})^2&=(\cA{3x})^2+2\cdot\cA{3x}\cdot\cB{2}+\cB{2}^2=\boxed{9x^2+12x+4} && \text{\nota{$a=3x$,\ $b=2$}}\\
(\cA{x^2}-\cB{5})^2&=(\cA{x^2})^2-2\cdot\cA{x^2}\cdot\cB{5}+\cB{5}^2=\boxed{x^4-10x^2+25}\\
(\cA{2x}+\cB{3y})(\cA{2x}-\cB{3y})&=(\cA{2x})^2-(\cB{3y})^2=\boxed{4x^2-9y^2}
\end{align*}
\end{exemple}

\vfill
\begin{minipage}[t]{0.5\linewidth}
\begin{compte}[Errores típicos]
\mal{(x+3)^2=x^2+9}\ \ {\small ¡falta el doble!}\\[2pt]
\be{(x+3)^2=x^2+6x+9}\\[5pt]
\mal{(3x)^2=3x^2}\ \ {\small el 3 también se eleva}\\[2pt]
\be{(3x)^2=9x^2}
\end{compte}
\end{minipage}\hfill
\begin{minipage}[t]{0.47\linewidth}
\begin{truc}[También al revés]
Si ves dos cuadrados y el doble producto\dots\\[4pt]
$x^2+14x+49=(x+7)^2$\\[2pt]
$x^2-16=(x+4)(x-4)$
\end{truc}
\end{minipage}

\vfill
% ============================================================ 3.5
\clearpage
\seccio{sec5}{División de polinomios}{Se hace como la división de números\dots{} pero con $x$.}

\vfill
\begin{minipage}[c]{0.34\linewidth}\centering
{\Large$\begin{array}{r|l}\cA{17} & \cB{5}\\ \cline{2-2} \cR{2} & \cK{3}\end{array}$}\\[6pt]
{\Large$\cA{17}=\cB{5}\cdot\cK{3}+\cR{2}$}
\end{minipage}\hfill
\begin{minipage}[c]{0.05\linewidth}\centering{\Huge\color{fmmuted}\faArrowRight}\end{minipage}\hfill
\begin{minipage}[c]{0.57\linewidth}
\begin{idea}
\centering{\Large$\cA{D(x)}=\cB{d(x)}\cdot\cK{C(x)}+\cR{R(x)}$}\\[4pt]
\raggedright\small\cA{\textbf{dividendo}} \ = \ \cB{\textbf{divisor}} $\cdot$ \cK{\textbf{cociente}} + \cR{\textbf{resto}}\\[2pt]
Acabamos cuando el grado del \cR{resto} es \textbf{menor} que el del \cB{divisor}. Si el resto es 0, la división es \textbf{exacta}.
\end{idea}
\end{minipage}

\vfill
\subsection{Paso a paso}
\begin{minipage}[c]{0.4\linewidth}
\pas{1} \textbf{Ordenamos} los dos polinomios.\\[5pt]
\pas{2} \textbf{Dividimos} el primer término del dividendo entre el primero del divisor.\\[5pt]
\pas{3} \textbf{Multiplicamos} por todo el divisor y lo escribimos debajo \textbf{cambiado de signo}. \textbf{Sumamos}.\\[5pt]
\pas{4} \textbf{Repetimos} hasta que no se pueda más.
\end{minipage}\hfill
\begin{minipage}[c]{0.57\linewidth}
\begin{exemple}[Ejemplo: \ $(x^3-2x^2+4x-5):(x-2)$]
\renewcommand{\arraystretch}{1.35}\large
\[\begin{array}{rrrr|l@{\qquad}l}
x^3 & -2x^2 & +4x & -5 & \,\cB{x-2} & \\ \cline{5-5}
\cR{-x^3} & \cR{+2x^2} & & & \,\cK{x^2+4} & \text{\nota{\pas{2}\ $x^3:x=\cK{x^2}$}}\\ \cline{1-2}
 & 0 & +4x & -5 & & \text{\nota{\pas{3}\ sumamos}}\\
 & & \cR{-4x} & \cR{+8} & & \text{\nota{\pas{4}\ $4x:x=\cK{+4}$}}\\ \cline{3-4}
 & & & \cR{\boxed{3}} & & \text{\nota{resto}}
\end{array}\]
\normalsize $\cK{C(x)=x^2+4}$\qquad $\cR{R=3}$
\end{exemple}
\end{minipage}

\vfill
\begin{minipage}[t]{0.52\linewidth}
\begin{compte}[¿Falta algún grado?]
Deja un \textbf{hueco} (o pon un 0):\\[3pt]
$x^4\qquad\ \ -3x^2+2x$\ \ \ \notab{faltan $x^3$ y el número}
\end{compte}
\end{minipage}\hfill
\begin{minipage}[t]{0.45\linewidth}
\begin{truc}[Compruébalo]
$\cB{(x-2)}\cdot\cK{(x^2+4)}+\cR{3}$\\
$=x^3-2x^2+4x-8+3$\\
$=x^3-2x^2+4x-5$ \ \be{}
\end{truc}
\end{minipage}

\vfill
% ============================================================ 3.6
\clearpage
\seccio{sec6}{Regla de Ruffini}{Un atajo para dividir entre \,$(x-a)$.}

\vfill
\begin{idea}
Cuando el divisor es \fmkey{$x-a$}, podemos dividir \textbf{mucho más rápido}: trabajamos \textbf{solo con los coeficientes}.
\end{idea}

\vfill
{\large Ejemplo: \ $(2x^3-5x^2+4):(x-3)$}

\vfill
\begin{minipage}[c]{0.44\linewidth}\centering
\ruffini{3}{2,-5,0,4}{6,3,9}{2,1,3,13}{4}\\[4pt]
\tcbox[colback=fmK!10,colframe=fmK,arc=3mm,boxrule=1.2pt,left=6pt,right=6pt]{\large$C(x)=2x^2+x+3$\qquad\cR{$R=13$}}\par
{\comic\small\color{fmmuted}el cociente tiene un grado menos}
\end{minipage}\hfill
\begin{minipage}[c]{0.53\linewidth}
\pas{1} Escribimos los \textbf{coeficientes} (con \textbf{0} donde falta un término). A la izquierda, el número \cB{$a$}.\\[5pt]
\pas{2} {\color{fmK}\faArrowDown}\ \textbf{Bajamos} el primer coeficiente.\\[5pt]
\pas{3} {\color{fmB}\faArrowRight}\ \textbf{Multiplicamos} por \cB{$a$} y lo ponemos en la columna siguiente.\\[5pt]
\pas{4} {\color{fmA}\faArrowDown}\ \textbf{Sumamos}. Y volvemos a \pas{3}.\\[5pt]
\pas{5} El \cR{\textbf{último}} es el \cR{\textbf{resto}}; los demás, el cociente.
\end{minipage}

\vfill
\begin{compte}[¡Cuidado con el signo de $a$!]
\begin{tabular}{@{}l@{\qquad}l@{}}
$x\cB{-3}$ \ $\to$ \ ponemos \ \cB{\textbf{3}} & $x\cB{+2}=x-(-2)$ \ $\to$ \ ponemos \ \cB{$\mathbf{-2}$}
\end{tabular}
\end{compte}

\vfill
\subsection{Teorema del resto}
\begin{minipage}[c]{0.5\linewidth}
\begin{idea}
\centering{\Large El resto de \ $P(x):(x-a)$ \ es \ $\boxed{R=P(a)}$}\\[4pt]
\raggedright\small $P(3)=2\cdot27-5\cdot9+4=13$ \ \be{} (¡es el mismo resto!)
\end{idea}
\end{minipage}\hfill
\begin{minipage}[c]{0.47\linewidth}
\begin{truc}[Si el resto es 0\dots]
$P(a)=0$ \ $\Rightarrow$ \ la división es \textbf{exacta}:
\begin{itemize}[leftmargin=1.1em,itemsep=0pt]
\item $a$ es una \textbf{raíz} de $P$,
\item $(x-a)$ es un \textbf{factor}: \ $P(x)=(x-a)\cdot C(x)$.
\end{itemize}
\end{truc}
\end{minipage}

\vfill
\begin{exemple}[Ejemplo: \ $(x^3+2x^2-5x-6):(x-2)$]
\begin{minipage}[c]{0.42\linewidth}\centering
\ruffini{2}{1,2,-5,-6}{2,8,6}{1,4,3,0}{4}
\end{minipage}\hfill
\begin{minipage}[c]{0.55\linewidth}
$C(x)=x^2+4x+3$,\quad \cR{$R=0$}\ \ \nota{¡exacta!}\\[6pt]
{\large$x^3+2x^2-5x-6=(x-2)(x^2+4x+3)$}\\[4pt]
\nota{2 es una raíz y $(x-2)$ es un factor}
\end{minipage}
\end{exemple}

\vfill
\clearpage
\subsection{Valor numérico y teorema del resto}

\vfill
\begin{minipage}[c]{0.5\linewidth}
\begin{idea}[Valor numérico]
$P(a)$: sustituimos $x$ por $a$.\\[3pt]
$P(x)=x^3-2x+5$: \ $P(-2)=-8+4+5=\boxed{1}$
\end{idea}
\end{minipage}\hfill
\begin{minipage}[c]{0.47\linewidth}
\begin{tcolorbox}[enhanced,colback=fmyellow!18,colframe=fmgold,arc=4mm,boxrule=1.6pt,before upper=\centering,top=3mm,bottom=3mm]
{\Large Teorema del resto}\par\smallskip
el resto de $P(x):(x-a)$ es\par\smallskip
{\fontsize{22}{26}\selectfont$r=P(a)$}
\end{tcolorbox}
\end{minipage}

\vfill
\begin{truc}[¿Por qué?]
$P(x)=(x-a)\cdot c(x)+r$. \ Si ponemos $x=a$, el primer sumando vale 0: \ $P(a)=r$.\\[2pt]
Dos maneras de calcular $P(a)$: sustituyendo, o haciendo Ruffini y mirando el resto (más rápido con grados altos).
\end{truc}

\vfill
\begin{tcbraster}[raster columns=2,raster equal height,raster column skip=10pt]
\begin{exemple}[Ejemplo: \ resto de $(x^5-3x^2+7):(x+1)$]
Sin dividir: \ $r=P(-1)=-1-3+7=\boxed{3}$
\end{exemple}
\begin{exemple}[Ejemplo: \ $k$ para que $(x^3+kx^2-5x+6):(x-3)$ sea exacta]
$P(3)=0$: \ $27+9k-15+6=0\ \to\ \boxed{k=-2}$
\end{exemple}
\end{tcbraster}

\vfill
\begin{exemple}[Ejemplo: \ $P(3)$ con Ruffini, para $P(x)=x^4-3x^3+2x-5$]
\begin{minipage}[c]{0.5\linewidth}\centering
\scalebox{0.78}{\ruffini{3}{1,-3,0,2,-5}{3,0,0,6}{1,0,0,2,1}{5}}
\end{minipage}\hfill
\begin{minipage}[c]{0.47\linewidth}
No hay término en $x^2$: ponemos un \textbf{0}.\\[4pt]
El resto es el valor numérico: \ $P(3)=\boxed{1}$\\[4pt]
Comprobamos: \ $81-81+6-5=1$ \ \be{}
\end{minipage}
\end{exemple}

\vfill
\begin{compte}[El signo de $a$]
Dividir entre $x+2$ es dividir entre $x-(-2)$: el resto es \be{P(-2)}, no \mal{P(2)}.
\end{compte}

\vfill
\begin{exemple}[Ejemplo: \ $a$ y $b$ para que $x^3+ax^2+bx-6$ sea divisible por $x-1$ y, al dividirlo por $x+1$, el resto sea $-4$]
$P(1)=0$: \ $1+a+b-6=0\ \to\ a+b=5$ \qquad $P(-1)=-4$: \ $-1+a-b-6=-4\ \to\ a-b=3$\\[4pt]
Resolvemos el sistema: \ $2a=8$ \ $\to$ \ \boxed{a=4,\ b=1}
\end{exemple}

\vfill
\begin{exemple}[Ejemplos]
$P(x)=2x^3-x^2+3$: \ $P(0)=3$, \ $P(1)=4$, \ $P(-1)=0$, \ $P(2)=15$\\[4pt]
Resto de $(x^3-mx+4):(x+2)$ igual a 10: \ $P(-2)=-8+2m+4=10\ \to\ \boxed{m=7}$\\[4pt]
¿Es exacta $(x^4-16):(x+2)$? \ $P(-2)=16-16=0$ \ \be{} sí: \ $x+2$ es un factor de $x^4-16$
\end{exemple}

\vfill

\vfill
% ============================================================ 3.7
\clearpage
\seccio{sec7}{Raíces de un polinomio}{Los valores que anulan el polinomio.}

\vfill
\begin{minipage}[c]{0.52\linewidth}
\begin{idea}[Raíz y teorema del factor]
$a$ es una \textbf{raíz} de $P(x)$ si $P(a)=0$.\\[4pt]
{\centering$P(a)=0\iff(x-a)$ es un \textbf{factor} de $P(x)$\par}\smallskip
Grado $n$: $n$ raíces como mucho. \ $x^2+1$: ninguna.
\end{idea}
\end{minipage}\hfill
\begin{minipage}[c]{0.45\linewidth}
\begin{tcolorbox}[fmbase,colback=fmK!8,colframe=fmK,coltitle=white,title={\faSearch\ \ Raíces enteras}]
\pas{1} Candidatos: divisores del \textbf{término indep.}\\[2pt]
\pas{2} Probamos: si $P(a)=0$, es raíz.\\[2pt]
\pas{3} Sin término indep.: factor común $x$.
\end{tcolorbox}
\end{minipage}

\vfill
\begin{exemple}[Ejemplo: \ raíces enteras de $P(x)=x^3-2x^2-5x+6$]
Candidatos: divisores de 6: \ $\pm1,\ \pm2,\ \pm3,\ \pm6$\\[4pt]
$P(1)=0$ \be{} \qquad $P(-1)=8$ \qquad $P(2)=-4$ \qquad $P(-2)=0$ \be{} \qquad $P(3)=0$ \be{}\\[4pt]
Grado 3 y ya tenemos 3: \ \boxed{x=1,\ x=-2,\ x=3}
\end{exemple}

\vfill
\begin{minipage}[c]{0.44\linewidth}\centering
\begin{tikzpicture}[x=0.7cm,y=0.7cm]
  \draw[fmmuted!30,line width=0.4pt] (-3,-1.5) grid (4,2.5);
  \draw[-{Latex[length=2mm]},fmmuted,line width=0.8pt] (-3,0) -- (4.2,0) node[right,font=\scriptsize] {$x$};
  \draw[-{Latex[length=2mm]},fmmuted,line width=0.8pt] (0,-1.5) -- (0,2.7) node[above,font=\scriptsize] {$y$};
  \draw[fmA,line width=1.6pt,domain=-2.3:3.3,samples=80] plot (\x,{0.25*(\x*\x*\x-2*\x*\x-5*\x+6)});
  \foreach \r in {-2,1,3} {\fill[fmR] (\r,0) circle (3pt); \node[below left=-1pt,font=\scriptsize\heavy,text=fmR] at (\r,0) {$\r$};}
\end{tikzpicture}\\
{\comic\small\color{fmmuted}$y=x^3-2x^2-5x+6$}
\end{minipage}\hfill
\begin{minipage}[c]{0.53\linewidth}
\begin{truc}[Raíces y gráfica]
Las raíces son los puntos donde la gráfica \textbf{corta el eje $X$}: allí $y=P(x)=0$.\\[3pt]
Una parábola, $x^2+1$, que no corta el eje no tiene raíces.
\end{truc}
\end{minipage}

\vfill
\begin{minipage}[t]{0.5\linewidth}
\begin{exemple}[Ejemplo: \ $x^3+x^2-4x-4$]
Candidatos $\pm1,\pm2,\pm4$: \ $\boxed{x=-1,\ 2,\ -2}$
\end{exemple}
\end{minipage}\hfill
\begin{minipage}[t]{0.47\linewidth}
\begin{compte}[¿Puede ser raíz el 4?]
De $x^3+2x^2-x+6$: \textbf{no}, porque 4 no es divisor de 6.
\end{compte}
\end{minipage}

\vfill
\begin{exemple}[Ejemplo: \ raíces de $x^3-3x+2$]
\begin{minipage}[c]{0.55\linewidth}
Candidatos $\pm1,\pm2$: \ $P(1)=0$ \be{}, \ $P(-1)=4$, \ $P(2)=4$, \ $P(-2)=0$ \be{}\\[4pt]
¿Solo dos raíces y el grado es 3? \ Una se \textbf{repite}: \\[2pt]
$x^3-3x+2=(x-1)^2(x+2)$ \ \nota{$x=1$ es raíz doble}
\end{minipage}\hfill
\begin{minipage}[c]{0.42\linewidth}\centering
\scalebox{0.7}{\ruffini{1}{1,0,-3,2}{1,1,-2}{1,1,-2,0}{4}}
\end{minipage}
\end{exemple}

\vfill
\begin{exemple}[Ejemplos]
\begin{minipage}[t]{0.48\linewidth}
$x^4-x^3-2x^2=x^2(x^2-x-2)$:\\ raíces \boxed{x=0\text{ (doble)},\ 2,\ -1} \ \nota{factor común}\\[4pt]
$2x^2-x-1$: \ $x=1$ y $x=-\frac12$ \ \nota{no todas son enteras}
\end{minipage}\hfill
\begin{minipage}[t]{0.49\linewidth}
Grado 3, raíces $1,\,-2,\,3$ y $P(0)=12$:\\
$P(x)=a(x-1)(x+2)(x-3)$, \ $P(0)=6a=12$\\
$\boxed{P(x)=2(x-1)(x+2)(x-3)}$
\end{minipage}
\end{exemple}

\vfill

\vfill
% ============================================================ 3.8
\clearpage
\seccio{sec8}{Factorización de polinomios}{Producto de factores del menor grado posible.}

\vfill
\begin{tcolorbox}[enhanced,colback=fmyellow!18,colframe=fmgold,arc=4mm,boxrule=1.6pt,before upper=\centering,top=2mm,bottom=2mm]
{\Large$P(x)=a_n\,(x-x_1)(x-x_2)\cdots(x-x_n)$}\par\smallskip
{\small $x_1,\dots,x_n$ son las raíces; si un factor se repite, la raíz es \textbf{doble}, \textbf{triple}\dots}
\end{tcolorbox}

\vfill
\begin{tcbraster}[raster columns=4,raster equal height,raster column skip=6pt]
\begin{tcolorbox}[fmbase,colback=fmdefbg,colframe=fmdef,coltitle=white,title={\pas{1}\ Factor común},left=2mm,right=2mm]
$2x^4-6x^3$\\ $=2x^3(x-3)$
\end{tcolorbox}
\begin{tcolorbox}[fmbase,colback=fmdefbg,colframe=fmdef,coltitle=white,title={\pas{2}\ Identidades},left=2mm,right=2mm]
$x^2-9=(x+3)(x-3)$\\ $x^2+6x+9=(x+3)^2$
\end{tcolorbox}
\begin{tcolorbox}[fmbase,colback=fmdefbg,colframe=fmdef,coltitle=white,title={\pas{3}\ Grado 2},left=2mm,right=2mm]
Ecuación:\\ $a(x-x_1)(x-x_2)$
\end{tcolorbox}
\begin{tcolorbox}[fmbase,colback=fmdefbg,colframe=fmdef,coltitle=white,title={\pas{4}\ Grado mayor},left=2mm,right=2mm]
Ruffini hasta llegar a grado 2.
\end{tcolorbox}
\end{tcbraster}

\vfill
\begin{exemple}[Ejemplo: \ $P(x)=x^4-x^3-7x^2+x+6$]
\begin{minipage}[c]{0.45\linewidth}\centering
\ruffini{1}{1,-1,-7,1,6}{1,0,-7,-6}{1,0,-7,-6,0}{5}
\end{minipage}\hfill
\begin{minipage}[c]{0.52\linewidth}
$x=1$: resto 0. Con el cociente y $x=-1$: \ $x^3-7x-6=(x+1)(x^2-x-6)$\\[3pt]
$x^2-x-6=0\ \to\ x=3,\ x=-2$\\[3pt]
$\boxed{P(x)=(x-1)(x+1)(x-3)(x+2)}$
\end{minipage}
\end{exemple}

\vfill
\begin{exemple}[Ejemplo: \ $P(x)=2x^3-3x^2-3x+2$ \ (coeficiente principal 2)]
\begin{minipage}[c]{0.4\linewidth}\centering
\scalebox{0.72}{\ruffini{-1}{2,-3,-3,2}{-2,5,-2}{2,-5,2,0}{4}}
\end{minipage}\hfill
\begin{minipage}[c]{0.57\linewidth}
$P(-1)=0$: \ $P(x)=(x+1)(2x^2-5x+2)$\\[3pt]
$2x^2-5x+2=0\ \to\ x=\dfrac{5\pm3}{4}=2,\ \dfrac12$\\[3pt]
$P(x)=2(x+1)(x-2)\left(x-\tfrac12\right)=\boxed{(x+1)(x-2)(2x-1)}$\\[2pt]
\notab{el 2 entra en el último paréntesis: \ $2\left(x-\frac12\right)=2x-1$}
\end{minipage}
\end{exemple}

\vfill
\begin{exemple}[Ejemplo: \ $P(x)=x^3-x^2+x-6$ \ (queda un factor de grado 2 sin raíces)]
\begin{minipage}[c]{0.4\linewidth}\centering
\scalebox{0.72}{\ruffini{2}{1,-1,1,-6}{2,2,6}{1,1,3,0}{4}}
\end{minipage}\hfill
\begin{minipage}[c]{0.57\linewidth}
$P(2)=0$: \ $P(x)=(x-2)(x^2+x+3)$\\[3pt]
$x^2+x+3=0$: \ $\Delta=1-12<0$, \ \textbf{no tiene solución}\\[3pt]
\raggedright Se queda tal cual en la factorización:\\[3pt] $\boxed{P(x)=(x-2)(x^2+x+3)}$
\end{minipage}
\end{exemple}

\vfill
\begin{exemple}[Ejemplos]
$x^3-6x^2+9x=x(x^2-6x+9)=\boxed{x(x-3)^2}$ \ \nota{$x=3$ doble}\\[4pt]
$x^2-16=(x+4)(x-4)$ \qquad $3x^2-12x+12=3(x-2)^2$ \qquad $x^3-x=x(x+1)(x-1)$\\[4pt]
$x^4-5x^2+4=(x-1)(x+1)(x-2)(x+2)$ \qquad $x^3+x^2-5x+3=(x-1)^2(x+3)$
\end{exemple}

\vfill

\vfill
% ============================================================ 3.9
\clearpage
\seccio{sec9}{Aplicación: fracciones algebraicas}{Cocientes de polinomios: simplificamos y operamos como con fracciones.}

\vfill
\begin{minipage}[t]{0.5\linewidth}
\begin{idea}[Fracción algebraica]
$\dfrac{P(x)}{Q(x)}$: solo tiene sentido si $Q(x)\neq0$.\\[3pt]
\textbf{Simplificar}: factorizamos y eliminamos los factores comunes.\\[3pt]
$\dfrac{x^3-4x}{x^3-4x^2+4x}=\dfrac{x(x+2)(x-2)}{x(x-2)^2}=\boxed{\dfrac{x+2}{x-2}}$
\end{idea}
\end{minipage}\hfill
\begin{minipage}[t]{0.47\linewidth}
\begin{compte}[Factores, no sumandos]
Solo se simplifican \textbf{factores}:\\[3pt]
\mal{\dfrac{x+3}{x}=3}\qquad \mal{\dfrac{x^2+1}{x^2}=1}
\end{compte}
\end{minipage}

\vfill
\begin{tcbraster}[raster columns=3,raster equal height,raster column skip=8pt]
\begin{tcolorbox}[fmbase,colback=fmdefbg,colframe=fmdef,coltitle=white,title={Suma y resta}]
Común denominador: el \textbf{m.c.m.} de los denominadores.
\end{tcolorbox}
\begin{tcolorbox}[fmbase,colback=fmdefbg,colframe=fmdef,coltitle=white,title={Producto},before upper=\centering]
$\dfrac PQ\cdot\dfrac RS=\dfrac{P\cdot R}{Q\cdot S}$
\end{tcolorbox}
\begin{tcolorbox}[fmbase,colback=fmdefbg,colframe=fmdef,coltitle=white,title={Cociente},before upper=\centering]
$\dfrac PQ:\dfrac RS=\dfrac{P\cdot S}{Q\cdot R}$
\end{tcolorbox}
\end{tcbraster}

\vfill
\begin{exemple}[Ejemplo: \ $\dfrac{x}{x-1}+\dfrac{2}{x^2-1}$]
$x^2-1=(x+1)(x-1)$: \ m.c.m. $=(x+1)(x-1)$ \qquad $=\dfrac{x(x+1)+2}{(x+1)(x-1)}=\boxed{\dfrac{x^2+x+2}{x^2-1}}$
\end{exemple}

\vfill
\begin{exemple}[Ejemplo: \ producto y cociente]
$\dfrac{x^2-9}{x^2+2x}\cdot\dfrac{x+2}{x-3}=\dfrac{(x+3)(x-3)(x+2)}{x(x+2)(x-3)}=\boxed{\dfrac{x+3}{x}}$ \qquad
$\dfrac{2x-4}{x^2}:\dfrac{x^2-4}{x}=\dfrac{2(x-2)\cdot x}{x^2(x+2)(x-2)}=\boxed{\dfrac{2}{x(x+2)}}$
\end{exemple}

\vfill
\begin{exemple}[Ejemplo: \ $\left(\dfrac1x-\dfrac{1}{x+1}\right):\dfrac{1}{x^2+x}$]
Primero el paréntesis: \ $\dfrac{(x+1)-x}{x(x+1)}=\dfrac{1}{x(x+1)}$ \qquad Después: \ $\dfrac{1}{x(x+1)}:\dfrac{1}{x(x+1)}=\boxed{1}$ \ \nota{si $x\neq0,\,-1$}
\end{exemple}

\vfill
\begin{truc}[Factoriza antes de operar]
Factoriza numeradores y denominadores \textbf{antes} de multiplicar: así simplificas enseguida.
\end{truc}

\vfill

\vfill
% ============================================================ resum
\resumtitol
\begin{tcbraster}[raster columns=2,raster equal height=rows,raster column skip=14pt,raster row skip=26pt,raster before skip=16pt]
\rtile{3.1}{Monomios}{Coeficiente y parte literal. Solo sumamos los \textbf{semejantes}.}
\rtile{3.2}{Polinomios}{Grado, coeficiente principal, término independiente. $P(a)$: sustituimos.}
\rtile{3.3}{Suma y resta}{Agrupamos los términos semejantes.\\[1ex] Restar: cambiamos todos los signos.}
\rtile{3.4}{Producto}{Cada término por cada término.\\[1ex] $(a\pm b)^2=a^2\pm2ab+b^2$,\\[1ex] $(a+b)(a-b)=a^2-b^2$.}
\rtile{3.5}{División}{$P=d\cdot c+r$, \ grado de $r$ menor que el de $d$.}
\rtile{3.6}{Ruffini}{División entre $x-a$. El resto es $P(a)$.}
\rtile{3.7}{Raíces}{$P(a)=0$.\\[1ex] Candidatos: divisores del término independiente.}
\rtile{3.8}{Factorizar}{Factor común, identidades, Ruffini y ecuación de grado 2.}
\rtile{3.9}{Fracciones}{Factorizamos y simplificamos factores, no sumandos.}
\end{tcbraster}
\end{document}
